% Authoritative exercise source. Preserve all derivations and diagrams.

\ExerciseStart{G3.13}{Eseguire una ricerca fra quattro alternative}{Collegamento}

\BookHeading{Problema, dati e obiettivo}

Cercare l’elemento 11 fra quattro alternative con Grover. Partire dalla sovrapposizione uniforme, marcare 11 e applicare il diffusore \(D=2\ket s\bra s-I\). Calcolare anche una seconda iterazione per osservare la sovrarotazione.

\ExTheory{\ExRef{c3-grover-primo}{3.12}}{L’oracolo cambia il segno dell’ampiezza marcata. Il diffusore trasforma ogni ampiezza reale \(a_j\) in \(2\bar a-a_j\), dove \(\bar a\) è la media dopo la marcatura. Nel piano \((\ket b,\ket g)\), \(\ket g=\ket{11}\) e \(\ket b=(\ket{00}+\ket{01}+\ket{10})/\sqrt3\).}

\ExPrediction{La preparazione ha coordinate \((\sqrt3/2,1/2)\), quindi forma \(30^\circ\) con l’asse delle alternative non marcate. Le due riflessioni dovrebbero portarla sull’asse marcato dopo una iterazione.}

\ExFigure{\begin{tikzpicture}[>=Stealth,exlabel]
\draw[exaxis] (-0.35,0)--(1.55,0) node[right] {$b$};\draw[exaxis] (0,-0.35)--(0,1.05) node[above] {$g$};\draw[->,Accent,line width=1.1pt] (0,0)--(0.866025,0.5) node[anchor=south west] {$ s $};
\end{tikzpicture}}{Piano degli stati collettivi ortonormali: asse orizzontale b, asse verticale g.}


\begin{ExTasks}

\item Scrivere la preparazione uniforme e le coordinate collettive.

\item Applicare la marcatura.

\item Calcolare la media delle quattro ampiezze.

\item Applicare la riflessione rispetto alla media a ogni componente.

\item Ripetere marcatura e diffusione una seconda volta e interpretare il risultato.

\end{ExTasks}

\ExSolution

\ExStep{Preparare il piano di ricerca}

\ExEq{\ket s&=\tfrac12(\ket{00}+\ket{01}+\ket{10}+\ket{11})\\
&=\tfrac{\sqrt3}{2}\ket b+\tfrac12\ket g.}
Si ha \(\sin\theta=1/2\), quindi \(\theta=\pi/6\).

\ExStep{Marcare la soluzione}

\ExEq{O\ket s&=\tfrac12\colv{1\\1\\1\\-1}.}
Nel piano collettivo la coordinata buona cambia da \(1/2\) a \(-1/2\).

\ExStep{Calcolare la media}

\ExEq{\bar a&=(1/2+1/2+1/2-1/2)/4=1/4,\\2\bar a&=1/2.}

\ExStep{Riflettere ogni componente}

Per ciascuna delle tre componenti non marcate si ottiene \(1/2-1/2=0\). Per la componente marcata si ottiene \(1/2-(-1/2)=1\). Quindi \(DO\ket s=\ket{11}\) e la probabilità di successo è uno.

\ExStep{Eseguire una seconda iterazione}

L’oracolo porta \((0,0,0,1)^T\) in \((0,0,0,-1)^T\). La media è \(-1/4\), dunque \(2\bar a=-1/2\). Le tre componenti non marcate diventano \(-1/2\); quella marcata diventa \(-1/2-(-1)=1/2\). Le probabilità sono tutte \(1/4\).

\ExStep{Interpretare la rotazione}

Dopo k iterazioni l’angolo è \((2k+1)\theta\). Per \(k=1\), \(3\theta=\pi/2\), quindi la componente buona è 1. Per \(k=2\), \(5\theta=5\pi/6\), quindi è \(1/2\), con probabilità \(1/4\). Il vettore ha oltrepassato la direzione ottimale.

\par\noindent\begin{minipage}{\linewidth}
\ExFigure{\begin{tikzpicture}[>=Stealth,exlabel]
\draw[exaxis] (-1.4,0)--(1.6,0) node[right] {$b$};\draw[exaxis] (0,-.9)--(0,1.5) node[above] {$g$};\draw[Ink!20] (0,0) circle(1);\draw[exother] (0,0)--(.866,-.5) node[right] {$Os$};\draw[exvec] (0,0)--(0,1) node[above left] {$DOs$};\draw[Accent!50,->] (.866,-.5) arc(-30:90:1);
\end{tikzpicture}}{La marcatura riflette rispetto a b; la diffusione riflette rispetto alla direzione s. La loro composizione ruota di \(2\theta\).}


\ExInterpret{La norma resta uno in ogni passaggio, mentre la probabilità della soluzione varia con l’angolo. Aggiungere iterazioni oltre il punto utile può ridurre il successo.}

\ExLab{Nella sezione 3.12 scegliere «Soluzione 11» e verificare il passaggio dalle quattro ampiezze uguali alla soluzione certa. Nella sezione 4.05 confrontare la stessa interpretazione con registri più grandi e più iterazioni.}


\end{minipage}
